\documentclass{beamer}
\usetheme{Singapore}
\usepackage[utf8]{inputenc}
\usecolortheme{crane}
\usepackage{graphicx}
\usepackage{standalone}
\usepackage{tikz}
\usetikzlibrary{arrows}
\usetikzlibrary{decorations.markings}
\usetikzlibrary{calc}
\usetikzlibrary{shapes,snakes}
\usetikzlibrary{decorations.text}
\usepackage{amsmath}
\usepackage{amsfonts}
\usepackage{amsthm}
\usepackage{mathtools}
\usepackage{tcolorbox}
\usepackage{float}
\usepackage{bm}
\usepackage{minted}
\usepackage{booktabs}

\definecolor{lightblue}{RGB}{124,190,255}
\definecolor{darkgreen}{RGB}{24,145,0}
\definecolor{darkorange}{RGB}{220,110,0}
\definecolor{bg}{RGB}{249, 251, 223}

\setbeamertemplate{section in toc}[sections numbered]
\beamertemplatenavigationsymbolsempty
\setbeamerfont{caption}{size=\tiny}


\title{Dulliau ac Offerynnau ar gyfer Ymchwil Ailgynhyrchiadwy}
\author{\textcolor{darkorange}{@GeraintPalmer}}
\date{Cynhadledd Wyddonol 2018}
\titlegraphic{
  \begin{tikzpicture}{width=2.35cm}
    \node at (0, 0) {\includegraphics[height=1.4cm]{ccclogo}};
    \node at (2.5, 0) {\includegraphics[height=1.4cm]{../cflogo}};
    \node at (5.5, 0) {\includegraphics[height=1.1cm]{SSI_Logo}};
  \end{tikzpicture}%
}

\begin{document}

\frame{\titlepage}

\begin{frame}
Pam ailgynhyrchiadwy?
\end{frame}

\begin{frame}
Beth yw ymchwil?
\end{frame}

\begin{frame}
\frametitle{\hfill}
\tableofcontents
\end{frame}

\section[]{Osgoi llafr llaw}
\begin{frame}
\frametitle{\hfill}
\tableofcontents[currentsection,currentsubsection]
\end{frame}

\begin{frame}
Click, click, click
\end{frame}

\begin{frame}
Excel
\end{frame}

\begin{frame}
Sgriptio
\end{frame}

\section[]{Offerynnau agored}
\begin{frame}
\frametitle{\hfill}
\tableofcontents[currentsection,currentsubsection]
\end{frame}

\begin{frame}
Trwyddedau
\end{frame}

\begin{frame}
Ffynhonell agored
\end{frame}

\section[]{Ymarferion cod gorau}
\begin{frame}
\frametitle{\hfill}
\tableofcontents[currentsection,currentsubsection]
\end{frame}

\begin{frame}[fragile]
\scriptsize{
\begin{minted}{python}
a = 1071
b = 462
C = []

while a != 0 and b != 0:
  i = 0
  while a >= b:
    a -= b
    i += 1
  a, b = b, a
  C.append(i)
C.append(a)
C.append(b)

print(max(C))
\end{minted}
}
\end{frame}

\begin{frame}[fragile]
\scriptsize{
\begin{minted}{python}
mwyaf = 1071
lleiaf = 462
ffactorau = []

# Algorithm Euclid
while mwyaf != 0 and lleiaf != 0:
  lluosrif = 0
  while mwyaf >= lleiaf:
    mwyaf -= lleiaf
    lluosrif += 1
  mwyaf, lleiaf = lleiaf, mwyaf
  # cyfnewid enwau'r newidynnau
  ffactorau.append(lluosrif)
ffactorau.append(mwyaf)
ffactorau.append(lleiaf)

print(max(ffactorau))
\end{minted}
}
\end{frame}

\begin{frame}[fragile]
\scriptsize{
\begin{minted}{python}
def tynnu_lluosrifau(mwyaf, lleiaf):
    """
    Faint o weithiau mae un rhif yn mynd mewn i rhif arall
    """
    lluosrif = 0
    while mwyaf >= lleiaf:
        mwyaf -= lleiaf
        lluosrif += 1
    return mwyaf, lleiaf, lluosrif

def ffactor_mwyaf_cyffredin(mwyaf, lleiaf):
    """
    Algorithm Euclid y canfod ffactor mwyaf cyffredin dau rhif
    """
    ffactorau = []
    while mwyaf != 0 and lleiaf != 0:
        lleiaf, mwyaf, lluosrif = tynnu_lluosrifau(mwyaf, lleiaf)
        ffactorau.append(lluosrif)
    ffactorau.append(mwyaf)
    ffactorau.append(lleiaf)
    return max(ffactorau)

print(ffactor_mwyaf_cyffredin(1071, 462))
\end{minted}
}
\end{frame}

\begin{frame}[fragile]
\begin{minted}{python}
assert tynnu_lluosrifau(31, 10) == (1, 10, 3)

assert ffactor_mwyaf_cyffredin(1071, 462) == 21
\end{minted}
\end{frame}

\begin{frame}
  \begin{center}
    \vspace{0.8cm}
    \includestandalone[width=0.85\textwidth]{triongl-ymarferion-gorau}
  \end{center}
\end{frame}

\section[]{Rheolaeth fersiwn}
\begin{frame}
\frametitle{\hfill}
\tableofcontents[currentsection,currentsubsection]
\end{frame}

\begin{frame}
VC?
\end{frame}

\begin{frame}
Adnoddau Cymraeg i'w ddod.
\end{frame}
\end{document}
